\chapter{Central Banks and the Short Rate}\label{ch:m2:central-banks-and-the-short-rate}

At 14:00 New York time on 16 September 2026 the Federal Open Market
Committee published a statement of a few paragraphs raising its target for
the federal funds rate by a quarter point, to a range of 3.75 to 4 percent.
Nobody at the Federal Reserve bought or sold a bond to make it so. From the
next morning the Fed paid 3.90\% on every dollar banks kept in their accounts
with it, lent against Treasury collateral at 4.00\% to any eligible
counterparty that asked, and took cash from money-market funds at 3.75\%: three rates
it simply announced, and between which the market rates of the world's
largest currency then moved. The same day the European Central Bank's quarter
point took effect; two days later the Bank of Japan raised its target from
around 1\% to around 1.25\%. Every
interest rate in this book, the yield of a thirty-year bond, the price of a
swap, the forward points of a currency, is built on top of the overnight rate
that these few institutions set. This chapter explains how they set it, and
how the market turned it into the benchmark that replaced the interbank rates
of the previous forty years.

\section{Reserves and the central-bank balance sheet}

\begin{definition}[Reserves]\label{def:m2:central-banks-and-the-short-rate:reserves}
\emph{Reserves}\index{reserves} are the balances that banks, and some other
institutions admitted to it, hold in accounts at the central bank. They are
the asset in which payments between banks are finally settled: when a
customer of one bank pays a customer of another, reserves move from the first
bank's account to the second's.
\end{definition}

\omterm{def:m2:central-banks-and-the-short-rate:reserves}{Reserves} are a liability of the central bank, like the banknotes it issues and
the account it keeps for the government. On the other side of its balance
sheet sit the assets it bought or lent against: government bonds, mortgage
securities, loans to banks. \Cref{fig:m2:central-banks-and-the-short-rate:balance}
draws the Fed's two sides on one Wednesday.

\begin{proposition}[Only the central bank changes the total]\label{prop:m2:central-banks-and-the-short-rate:total}
Payments among banks and their customers move \omterm{def:m2:central-banks-and-the-short-rate:reserves}{reserves} between accounts but
do not change their total. The total changes only when an item of the central
bank's balance sheet changes: it rises when the central bank buys an asset or
lends, and falls when it sells or is repaid, or when another of its
liabilities (currency, the government's account, reverse repos) grows at the
expense of \omterm{def:m2:central-banks-and-the-short-rate:reserves}{reserves}.
\end{proposition}
\begin{proof}
The central bank's balance sheet satisfies assets $=$ \omterm{def:m2:central-banks-and-the-short-rate:reserves}{reserves} $+$ other
liabilities $+$ capital. A payment between two banks debits one reserve
account and credits another: no term of the identity moves. \omterm{def:m2:central-banks-and-the-short-rate:reserves}{Reserves} can
therefore change only if assets, capital or another liability change.
\end{proof}

\begin{example}[Tax day]\label{ex:m2:central-banks-and-the-short-rate:tax}
Companies pay USD~100 billion of taxes into the Treasury's account at the Fed.
Each payer's bank loses \omterm{def:m2:central-banks-and-the-short-rate:reserves}{reserves}; the Treasury's account gains the same
amount. The banking system has USD~100 billion fewer \omterm{def:m2:central-banks-and-the-short-rate:reserves}{reserves}, although
nobody decided anything about monetary policy. Movements in the government's
account are among the largest of these \emph{autonomous factors}, and a central
bank that wants a stable overnight rate must either offset them or hold so
many \omterm{def:m2:central-banks-and-the-short-rate:reserves}{reserves} that they do not matter.
\end{example}

\begin{omfigure}
\begin{tikzpicture}
\begin{axis}[width=11cm,height=5.2cm,ybar stacked,bar width=34pt,
  symbolic x coords={assets,liabilities},xtick=data,
  ylabel={USD billion},ymin=0,ymax=7200,enlarge x limits=0.55,
  tick label style={font=\small},label style={font=\small},
  legend columns=4,legend cell align=left,legend style={font=\footnotesize,at={(0.5,-0.2)},anchor=north,draw=none,fill=none,/tikz/every even column/.append style={column sep=6pt}},
  table/col sep=comma]
\addplot[fill=omDef!70,draw=omDef] table[x=side,y=treasuries]{figdata/markets-2/01-central-banks-and-the-short-rate/balance.csv};
\addplot[fill=omDef!35,draw=omDef] table[x=side,y=mbs]{figdata/markets-2/01-central-banks-and-the-short-rate/balance.csv};
\addplot[fill=gray!30,draw=gray] table[x=side,y=otherassets]{figdata/markets-2/01-central-banks-and-the-short-rate/balance.csv};
\addplot[fill=omThm!70,draw=omThm] table[x=side,y=reserves]{figdata/markets-2/01-central-banks-and-the-short-rate/balance.csv};
\addplot[fill=omProp!60,draw=omProp] table[x=side,y=tga]{figdata/markets-2/01-central-banks-and-the-short-rate/balance.csv};
\addplot[fill=omProp!25,draw=omProp] table[x=side,y=rrp]{figdata/markets-2/01-central-banks-and-the-short-rate/balance.csv};
\addplot[fill=omThm!25,draw=omThm] table[x=side,y=currency]{figdata/markets-2/01-central-banks-and-the-short-rate/balance.csv};
\legend{Treasuries,MBS,other assets,reserves,Treasury account,reverse repos,currency and rest}
\end{axis}
\end{tikzpicture}

\omcaption{The Federal Reserve's balance sheet on Wednesday 16 September 2026:
6\,747 billion dollars of assets, two thirds of them Treasuries, financed by
3\,014 billion of \omterm{def:m2:central-banks-and-the-short-rate:reserves}{reserves}, 877 billion in the Treasury's account, 324
billion of reverse repos, and currency, other liabilities and capital
together (the remainder, 2\,532 billion). Data: Federal Reserve, H.4.1
release of 17 September 2026.}\label{fig:m2:central-banks-and-the-short-rate:balance}
\end{omfigure}

\section{Corridor and floor systems}

\begin{definition}[Policy rate and standing facilities]\label{def:m2:central-banks-and-the-short-rate:policy}
A central bank's \emph{policy rate}\index{policy rate} is the interest rate
it announces as the stance of monetary policy and steers the market's
overnight rates towards. A \emph{standing facility}\index{standing facility}
is an operation available at its counterparties' initiative, in any amount
(against collateral where it lends), at a rate the central bank fixes: a
\emph{deposit facility} that pays interest on \omterm{def:m2:central-banks-and-the-short-rate:reserves}{reserves} and a \emph{lending
facility} that lends \omterm{def:m2:central-banks-and-the-short-rate:reserves}{reserves} overnight.
\end{definition}

\begin{proposition}[The facilities bound the overnight rate]\label{prop:m2:central-banks-and-the-short-rate:bounds}
Let $i_D < i_L$ be the rates of the deposit and lending facilities. A bank
with access to both never lends \omterm{def:m2:central-banks-and-the-short-rate:reserves}{reserves} overnight below $i_D$ and never
borrows above $i_L$ (against the same collateral). Among such banks the
overnight rate lies in $[i_D, i_L]$.
\end{proposition}
\begin{proof}
Lending in the market at $r < i_D$ earns less than leaving the \omterm{def:m2:central-banks-and-the-short-rate:reserves}{reserves} at
the deposit facility; borrowing at $r > i_L$ costs more than the lending
facility. Neither side accepts such a trade.
\end{proof}

\begin{definition}[Corridor and floor systems]\label{def:m2:central-banks-and-the-short-rate:corridor}
In a \emph{corridor system}\index{corridor system} the central bank keeps
\omterm{def:m2:central-banks-and-the-short-rate:reserves}{reserves} scarce, so that banks trade them among themselves and the overnight
rate settles inside $[i_D, i_L]$, near a target it steers by adjusting the
supply of \omterm{def:m2:central-banks-and-the-short-rate:reserves}{reserves} day by day. In a \emph{floor system}\index{floor system}
it supplies \omterm{def:m2:central-banks-and-the-short-rate:reserves}{reserves} beyond what banks need, so that the marginal reserve is
worth only the deposit rate, and the overnight rate sits at, or just below,
$i_D$.
\end{definition}

\Cref{fig:m2:central-banks-and-the-short-rate:corridor} is the whole theory in
one curve. Banks' demand for \omterm{def:m2:central-banks-and-the-short-rate:reserves}{reserves} falls with their price: when \omterm{def:m2:central-banks-and-the-short-rate:reserves}{reserves}
are scarce, a bank short at the end of the day pays up to the lending rate;
when they are abundant, no bank pays more than the deposit rate to hold one
more. Where the supply line crosses the curve is the overnight rate. On the
steep part, a small change in supply moves the rate, and the central bank
must forecast the autonomous factors of \cref{ex:m2:central-banks-and-the-short-rate:tax}
every morning. On the flat part it can ignore them: the rate is whatever it
pays on deposits, and moving that rate moves the market.

\begin{omfigure}
\begin{tikzpicture}
\begin{axis}[width=11cm,height=5cm,
  xlabel={reserves supplied (arbitrary units)},ylabel={overnight rate (\%)},
  tick label style={font=\small},label style={font=\small},
  xmin=0,xmax=6,ymin=2.4,ymax=3.0,grid=major,grid style={gray!25},table/col sep=comma]
\addplot[omDef,very thick] table[x=reserves,y=rate]{figdata/markets-2/01-central-banks-and-the-short-rate/corridor.csv};
\draw[omThm,thick,dashed] (axis cs:0,2.90) -- (axis cs:6,2.90);
\draw[omThm,thick,dashed] (axis cs:0,2.50) -- (axis cs:6,2.50);
\draw[omProp,thick] (axis cs:2.6,2.4) -- (axis cs:2.6,3.0);
\draw[omProp,thick] (axis cs:5.2,2.4) -- (axis cs:5.2,3.0);
\node[font=\footnotesize,anchor=south west] at (axis cs:0.1,2.905) {lending facility 2.90};
\node[font=\footnotesize,anchor=north west] at (axis cs:0.1,2.49) {deposit facility 2.50};
\node[font=\footnotesize,anchor=west,text=omProp] at (axis cs:2.65,2.78) {corridor};
\node[font=\footnotesize,anchor=east,text=omProp] at (axis cs:5.15,2.78) {floor};
\end{axis}
\end{tikzpicture}

\omcaption{A stylised demand curve for \omterm{def:m2:central-banks-and-the-short-rate:reserves}{reserves} between the ECB's facility
rates of 16 September 2026. With supply at the left vertical line (scarce
\omterm{def:m2:central-banks-and-the-short-rate:reserves}{reserves}) the rate is set on the steep part and moves with every change in
supply; at the right one (ample \omterm{def:m2:central-banks-and-the-short-rate:reserves}{reserves}) it rests on the deposit rate. The
curve's shape is illustrative; the two facility rates are the ECB's.
Data: the chapter's simulation.}\label{fig:m2:central-banks-and-the-short-rate:corridor}
\end{omfigure}

\begin{example}[The euro corridor]\label{ex:m2:central-banks-and-the-short-rate:ecb}
From 16 September 2026 the ECB's deposit facility pays 2.50\%, its weekly
main refinancing operations lend at 2.65\% and its marginal lending facility
at 2.90\%. The corridor is 40 basis points wide. Since 18 September 2024 the
refinancing rate sits only 15 basis points above the deposit rate (it had
been 50), a change announced in March 2024 so that banks would borrow in the
weekly operations as \omterm{def:m2:central-banks-and-the-short-rate:reserves}{reserves} shrink and the overnight rate would stay close
to the deposit rate. The euro area has, in effect, chosen a floor while
keeping the corridor as its outer bounds.
\end{example}

\begin{remark}[Below the floor]\label{rem:m2:central-banks-and-the-short-rate:below}
The bound of \cref{prop:m2:central-banks-and-the-short-rate:bounds} holds only
for institutions that can use the deposit facility. Cash-rich institutions
that cannot, such as money-market funds, lend to banks at whatever the banks
will pay, and the banks deposit the proceeds at the central bank for the
difference. That is why an overnight benchmark can print just below a floor:
on 22 September 2026 the euro short-term rate was 2.440\%, six basis points
under the deposit rate. The Fed's answer to the same leak is its reverse repo
facility, open to money-market funds and government-sponsored enterprises as
well as banks and dealers, which puts a second, lower floor under the market.
\end{remark}

\begin{dated}{2026-09}{Four central banks' administered rates}\label{dat:m2:central-banks-and-the-short-rate:rates}
\textbf{Federal Reserve} (from 17 September 2026): target range for the
federal funds rate 3.75--4.00\%; interest on reserve balances 3.90\%;
standing repo facility 4.00\%; overnight reverse repo facility 3.75\% (USD~160
billion per counterparty and day); primary credit 4.00\%.
\textbf{ECB} (from 16 September 2026): deposit facility 2.50\%, main
refinancing operations 2.65\%, marginal lending facility 2.90\%.
\textbf{Bank of England}: Bank Rate 3.75\%, held on 16 September 2026 by six
votes to three. \textbf{Bank of Japan} (from 24 September 2026): target for
the uncollateralised overnight call rate around 1.25\%; complementary
deposit facility 1.25\%; basic loan rate 1.50\%.
\end{dated}

\section{Open-market operations and balance-sheet policy}

\begin{definition}[Open-market operation]\label{def:m2:central-banks-and-the-short-rate:omo}
An \emph{open-market operation}\index{open-market operation} is a purchase or
sale of securities, or a repo or reverse repo, that the central bank
initiates, in an amount it chooses, to change the supply of \omterm{def:m2:central-banks-and-the-short-rate:reserves}{reserves}.
\end{definition}

\begin{definition}[Quantitative easing]\label{def:m2:central-banks-and-the-short-rate:qe}
\emph{Quantitative easing}\index{quantitative easing} is the purchase of
large amounts of longer-term securities, government bonds and sometimes
mortgage or corporate bonds, to lower longer-term yields once the \omterm{def:m2:central-banks-and-the-short-rate:policy}{policy rate}
can go no lower or no lower usefully. Its reversal, letting the securities
mature without replacing them or selling them, is \emph{quantitative
tightening}.
\end{definition}

\begin{example}[One purchase, three balance sheets]\label{ex:m2:central-banks-and-the-short-rate:purchase}
The central bank buys USD~10 billion of bonds from a pension fund. The
central bank: securities $+10$, \omterm{def:m2:central-banks-and-the-short-rate:reserves}{reserves} (of the fund's bank) $+10$. The
fund's bank: \omterm{def:m2:central-banks-and-the-short-rate:reserves}{reserves} $+10$, the fund's deposit $+10$. The fund: bonds
$-10$, deposit $+10$. The purchase created \omterm{def:m2:central-banks-and-the-short-rate:reserves}{reserves} and a bank deposit of
the same size; it did not create a loan, and the bank cannot ``lend out'' the
\omterm{def:m2:central-banks-and-the-short-rate:reserves}{reserves} to anyone but another bank, where they remain \omterm{def:m2:central-banks-and-the-short-rate:reserves}{reserves}
(\cref{prop:m2:central-banks-and-the-short-rate:total}).
\end{example}

In a \omterm{def:m2:central-banks-and-the-short-rate:corridor}{corridor system} \omterm{def:m2:central-banks-and-the-short-rate:omo}{open-market operations} \emph{are} monetary policy: the
desk adds or drains \omterm{def:m2:central-banks-and-the-short-rate:reserves}{reserves} each day to hold the rate at its target. In a
\omterm{def:m2:central-banks-and-the-short-rate:corridor}{floor system} they are not: the rate is set by the deposit rate, and the size
of the balance sheet becomes a separate instrument, used for longer-term
yields (\cref{def:m2:central-banks-and-the-short-rate:qe}) or for keeping
\omterm{def:m2:central-banks-and-the-short-rate:reserves}{reserves} safely on the flat part of the curve. That separation is the point
of a floor. It is also its cost: a central bank that pays interest on
trillions of \omterm{def:m2:central-banks-and-the-short-rate:reserves}{reserves} pays banks for holding them, in public.

The Fed's recent history shows both instruments at work. It shrank its
securities holdings by more than USD~2.2 trillion from June 2022, stopped on
1 December 2025 once it judged \omterm{def:m2:central-banks-and-the-short-rate:reserves}{reserves} close to the edge of ample, and in
December 2025 began buying Treasury bills again, not to ease policy but to
keep \omterm{def:m2:central-banks-and-the-short-rate:reserves}{reserves} growing with the economy. Its implementation note of September
2026, the same one that raised rates, repeats the instruction to buy bills as
needed: the rate went up and the balance sheet kept growing, without
contradiction.

\section{Overnight benchmarks and the end of the interbank offered rates}

\begin{definition}[Interbank offered rate]\label{def:m2:central-banks-and-the-short-rate:ibor}
An \emph{interbank offered rate}\index{interbank offered rate} (IBOR) is a
benchmark for unsecured term borrowing between banks, one to twelve months,
computed each day from submissions by a panel of banks of the rates at which
they could borrow. LIBOR, for five currencies, was the largest; trillions of
loans, bonds and derivatives referenced it.
\end{definition}

\begin{definition}[Overnight benchmark rate]\label{def:m2:central-banks-and-the-short-rate:rfr}
An \emph{overnight benchmark rate}\index{overnight benchmark rate} (also
called a risk-free rate, RFR) is a benchmark computed from actual overnight
transactions, published by a central bank or an administrator on the next
business day. It is either secured (repo against government bonds) or
unsecured (deposits taken by banks).
\end{definition}

The transition from the first kind to the second was forced by the sparseness
of the market the first measured: few banks still borrowed from each other
unsecured for three months, and a benchmark built on judgement rather than
transactions proved both fragile and manipulable. On 5 March 2021 the UK
regulator announced the end dates: most LIBOR settings after 31 December
2021, the main dollar settings after 30 June 2023. Synthetic dollar settings
were published as a bridge for legacy contracts until 30 September 2024;
since then all thirty-five LIBOR settings have ceased.

\begin{dated}{2026-09}{The four main overnight benchmarks}\label{dat:m2:central-banks-and-the-short-rate:benchmarks}
\textbf{SOFR} (US dollar, secured): overnight Treasury repo, tri-party,
general-collateral and cleared bilateral, with specials filtered out;
published by the Federal Reserve Bank of New York at about 08:00 New York
time for the previous business day. \textbf{\texteuro STR} (euro,
unsecured): euro-area banks' overnight wholesale borrowing; published by the
ECB each TARGET business day for the previous one (22 September 2026: 2.440\%
on EUR~69.9 billion, 929 transactions, 47 banks). \textbf{SONIA} (sterling,
unsecured): what banks pay to borrow sterling overnight from financial
institutions; published by the Bank of England at 09:00 London time the
next business day.
\textbf{Japan}: the uncollateralised overnight call rate, the Bank of Japan's
operating target.
\end{dated}

An overnight rate is known one day at a time, but a loan or a swap pays
interest for a quarter or a year. The market's answer is to compound the
daily rates over the period and pay at its end.

\begin{definition}[Compounding in arrears]\label{def:m2:central-banks-and-the-short-rate:arrears}
Over an interest period of $N$ calendar days with business days $d_1, \dots,
d_M$, fixings $r_1, \dots, r_M$ and weights $n_i$, the calendar days from
$d_i$ to the next business day, the rate \emph{compounded in
arrears}\index{compounding in arrears} is
\[
R \;=\; \Bigl[\,\prod_{i=1}^{M}\Bigl(1 + \frac{r_i\,n_i}{B}\Bigr) - 1\Bigr]\frac{B}{N},
\]
with $B = 360$ for the dollar and euro and 365 for sterling. It is known only
after the last fixing. Three conventions bring the payment date forward: a
\emph{lookback} of $p$ business days uses $r$ from $p$ days before each $d_i$
with the period's own weights; an \emph{observation shift} takes rates and
weights from a period moved back by $p$ business days; a \emph{lockout}
freezes the last $k$ fixings at the one before them.
\end{definition}

\begin{proposition}[Compounding adds a second-order term]\label{prop:m2:central-banks-and-the-short-rate:convexity}
Let $R_s = \sum_i r_i n_i / N$ be the day-weighted simple average. If all
$r_i \ge 0$ then $R \ge R_s$; for a constant rate $r$ fixed on every calendar
day,
\[
R - r \;=\; \frac{r^2 (N-1)}{2B} + O\!\left(\frac{r^3N^2}{B^2}\right).
\]
\end{proposition}
\begin{proof}
With $x_i = r_i n_i / B \ge 0$, $\prod (1 + x_i) - 1 = \sum x_i + \sum_{i<j}
x_i x_j + \dots \ge \sum x_i$, and $\sum x_i \cdot B/N = R_s$. For a constant
daily rate, $(1 + r/B)^N - 1 = Nr/B + \binom{N}{2}(r/B)^2 + O((Nr/B)^3)$;
multiply by $B/N$.
\end{proof}

\begin{example}[September 2026, illustrated]\label{ex:m2:central-banks-and-the-short-rate:month}
Take illustrative fixings of 3.62\% up to 16 September, 3.87\% from the
hike, and 3.95\% on the quarter-end day, 30 September
(\cref{fig:m2:central-banks-and-the-short-rate:month}). Over 1 September to
1 October (30 days, 21 fixings, Labor Day on 7 September) the day-weighted
average is 3.7393\% and the compounded rate 3.7448\%: compounding adds 0.54
basis points, against the 0.56 of \cref{prop:m2:central-banks-and-the-short-rate:convexity}
for a flat 3.74\%. With a five-day lookback the rate is 3.6836\%, 6.12 basis
points lower, USD~5\,098 on USD~100 million: the lookback moves five days of
the new rate into the next period, where the borrower will pay them.
\end{example}

\begin{omfigure}
\begin{tikzpicture}
\begin{axis}[width=11cm,height=5cm,
  xlabel={day of September 2026},ylabel={rate (\%)},
  tick label style={font=\small},label style={font=\small},
  xmin=1,xmax=31,ymin=3.55,ymax=4.0,grid=major,grid style={gray!25},
  legend columns=2,legend style={font=\footnotesize,at={(0.5,-0.3)},anchor=north,draw=none,fill=none,/tikz/every even column/.append style={column sep=8pt}},
  table/col sep=comma]
\addplot[omDef,very thick,const plot] table[x=day,y=fixing]{figdata/markets-2/01-central-banks-and-the-short-rate/month.csv};
\addplot[omThm,very thick,const plot,dashed] table[x=day,y=compounded]{figdata/markets-2/01-central-banks-and-the-short-rate/month.csv};
\legend{overnight fixing (held over weekends and the holiday),compounded from 1 September}
\end{axis}
\end{tikzpicture}

\omcaption{The illustrative fixings of \cref{ex:m2:central-banks-and-the-short-rate:month}
and the rate compounded from the start of the period to each day. The
compounded rate absorbs a change only gradually: each day at the new level is
one more weight in an average over all the days so far. Friday fixings weigh
three days, the Friday before Labor Day four. The
fixings are illustrative, not published SOFR. Data: the chapter's
tutorial.}\label{fig:m2:central-banks-and-the-short-rate:month}
\end{omfigure}

\begin{definition}[Fallback spread]\label{def:m2:central-banks-and-the-short-rate:fallback}
A \emph{fallback spread}\index{fallback spread} is the fixed adjustment added
to a compounded overnight rate when it replaces an \omterm{def:m2:central-banks-and-the-short-rate:ibor}{interbank offered rate} in
an existing contract, to compensate for the credit and term premium the old
rate contained. For LIBOR it was set as the five-year median of the
difference between the two rates, fixed on 5 March 2021: for dollar LIBOR
against SOFR, 11.448 basis points at one month, 26.161 at three months and
42.826 at six months.
\end{definition}

\begin{omfigure}
\begin{tikzpicture}[font=\small]
\draw[thick,-{Stealth[length=2.5mm]}] (0,0) -- (13,0) node[right]{time};
\draw[omDef,line width=4pt] (3.5,0.55) -- (11.5,0.55);
\node[above=3pt,text=omDef,font=\footnotesize] at (7.5,0.6) {interest period: weights $n_i$};
\draw[omThm,line width=4pt] (1.5,-0.55) -- (9.5,-0.55);
\node[below=3pt,text=omThm,font=\footnotesize] at (5.5,-0.6) {observation period: rates $r_i$ (and, with a shift, weights)};
\foreach \x/\t in {3.5/start,11.5/end}{\draw[thick] (\x,0.15) -- (\x,-0.15); \node[anchor=south,font=\footnotesize\bfseries] at (\x,0.8) {\t};}
\draw[-{Stealth[length=2mm]},omProp,thick] (3.5,-0.25) -- (1.6,-0.25);
\draw[-{Stealth[length=2mm]},omProp,thick] (11.5,-0.25) -- (9.6,-0.25);
\node[font=\footnotesize,text=omProp] at (10.5,-1.45) {$p$ business days};
\draw[decorate,decoration={brace,amplitude=4pt},thick] (9.5,1.55) -- (11.5,1.55) node[midway,above=5pt,font=\footnotesize,align=center] {last rate known: payment\\computed, sent on the end date};
\end{tikzpicture}

\omcaption{A lookback: each day of the interest period takes the fixing of
$p$ business days earlier, so the last rate needed is known $p$ days before
the period ends and the payment can be computed and sent on its end date.
With an observation shift the weights, too, come from the lower
period.}\label{fig:m2:central-banks-and-the-short-rate:lookback}
\end{omfigure}

\section{Tutorial: compounding a month in arrears}\label{tut:m2:central-banks-and-the-short-rate:compound}

\begin{tutorial}
\textbf{Goal.} Compound a month of overnight fixings on a real calendar,
with and without a lookback, and compare with the simple average.
\textbf{End state:} the four numbers of
\cref{ex:m2:central-banks-and-the-short-rate:month} and
\cref{fig:m2:central-banks-and-the-short-rate:month}.
\begin{tutsteps}
\item \textbf{Business days and weights.} A fixing accrues until the next
business day; a Friday's accrues three days.
\omcode{code/firm/rfr/firm_rfr.py}{32}{45}{Business days of a period and the calendar days each fixing accrues.}\label{lst:m2:central-banks-and-the-short-rate:weights}
You should see weights \texttt{[1, 4, 1, 1]} for 3 to 10 September 2026: the
Friday before Labor Day counts four days.
\item \textbf{Which fixing, which weight.} One function returns, for each
step of the period, the date of the fixing used and the days it accrues, for
all three conventions.
\omcode{code/firm/rfr/firm_rfr.py}{59}{72}{Fixing dates and weights under a lookback, an observation shift or a lockout.}\label{lst:m2:central-banks-and-the-short-rate:obs}
\item \textbf{Compound.}
\omcode{code/firm/rfr/firm_rfr.py}{75}{82}{The compounded rate of the period, annualised.}\label{lst:m2:central-banks-and-the-short-rate:compound}
\item \textbf{Run the month.} \texttt{short\_rate\_demo.month\_summary()}
prints 3.7393\% (simple), 3.7448\% (compounded) and 3.6836\% (five-day
lookback); \texttt{fig\_short\_rate.py} writes the chart's data.
\end{tutsteps}
\textbf{What to change next.} Replace the lookback by a two-day lockout, and
by a five-day observation shift, and explain why the second gives exactly the
lookback's number in this month (\cref{exo:m2:central-banks-and-the-short-rate:7}).
Then move the hike to the last week and watch the lookback's error grow.
\end{tutorial}

\section{Build: the overnight-rate compounding engine}\label{bld:m2:central-banks-and-the-short-rate:rfr}

\begin{build}
\bfield{Purpose} Every floating leg in the miniature firm, loans, swaps
(\cref{ch:m2:interest-rate-swaps}), futures on overnight rates
(\cref{ch:m2:short-term-interest-rate-futures}), repo accruals, pays a
compounded overnight rate. This component computes it, once and correctly.
\bfield{Interface} \texttt{Calendar(holidays)} with \texttt{is\_business\_day},
\texttt{add(d, n)}, \texttt{business\_days(start, end)};
\texttt{Convention(lookback, shift, lockout, basis)};
\texttt{observations(cal, start, end, conv)};
\texttt{compound\_in\_arrears(fixings, cal, start, end, conv)};
\texttt{simple\_average(\ldots)}; \texttt{rate\_from\_index(i0, i1, start,
end)}.
\bfield{Rules} Business days exclude weekends and the calendar's holidays;
weights run to the next business day and stop at the period end; rates are
decimals and may be negative; the basis is 360 unless the currency says 365.
\bfield{Acceptance tests} \texttt{code/firm/rfr/tests/}: a Friday counts three
days and the Friday before a Monday holiday four; a flat rate compounds to
the closed form and exceeds its simple average; a lookback keeps the
period's weights; an observation shift takes the shifted period's; a lockout
repeats the last free fixing; the published-index route gives the same rate;
negative rates.
\bfield{Stretch} Rounding as a published compounded index rounds (to eight
decimals); a payment-delay convention; a holiday calendar loaded from the
firm's reference data (One Quant Book 1, chapter 28).
\end{build}

\begin{omsources}
\item Board of Governors of the Federal Reserve System, \emph{Implementation
Note issued September 16, 2026}; statistical release H.4.1, 17 September
2026; minutes of the FOMC, December 2025.
\item European Central Bank, \emph{Key ECB interest rates}; press release
\emph{Changes to the operational framework for implementing monetary
policy}, 13 March 2024; \emph{Euro short-term rate}.
\item Bank of England, \emph{Monetary Policy Summary}, September 2026;
\emph{SONIA interest rate benchmark}. Bank of Japan, \emph{Change in the
Guideline for Money Market Operations}, 18 September 2026.
\item Federal Reserve Bank of New York, \emph{Secured Overnight Financing
Rate}; \emph{Reverse repo FAQ}.
\item Financial Conduct Authority, \emph{The end of LIBOR}, 2024; ISDA,
statement on the FCA announcement, 5 March 2021; ARRC recommended spread
adjustments, 2021.
\end{omsources}

\section{Exercises}

\begin{exercise}[$\star$]\label{exo:m2:central-banks-and-the-short-rate:1}
An interest period runs from Thursday to the following Tuesday, with fixings
of 3.90\% (Thursday), 3.88\% (Friday) and 3.92\% (Monday). Give the rate
\omterm{def:m2:central-banks-and-the-short-rate:arrears}{compounded in arrears} (basis 360) and the simple average.
\end{exercise}

\begin{exercise}[$\star$]\label{exo:m2:central-banks-and-the-short-rate:2}
With the ECB rates of \cref{dat:m2:central-banks-and-the-short-rate:rates},
give the width of the corridor and the refinancing rate's distance from the
deposit rate. The euro short-term rate prints at 2.440\%. Which bound of
\cref{prop:m2:central-banks-and-the-short-rate:bounds} does that seem to
violate, and why does it not?
\end{exercise}

\begin{exercise}[$\star$]\label{exo:m2:central-banks-and-the-short-rate:3}
The central bank sells USD~10 billion of bonds to an insurance company. Write
the changes on the balance sheets of the central bank, the insurer's bank and
the insurer.
\end{exercise}

\begin{exercise}[$\star\star$]\label{exo:m2:central-banks-and-the-short-rate:4}
Estimate the difference, in basis points, between a flat 4\% overnight rate
compounded daily over 91 calendar days and its simple average, first with
\cref{prop:m2:central-banks-and-the-short-rate:convexity} and then exactly.
\end{exercise}

\begin{exercise}[$\star\star$]\label{exo:m2:central-banks-and-the-short-rate:5}
Tax receipts of USD~150 billion flow into the Treasury's account on one day.
What happens to the overnight rate in a \omterm{def:m2:central-banks-and-the-short-rate:corridor}{corridor system} and in a \omterm{def:m2:central-banks-and-the-short-rate:corridor}{floor
system}? What does each central bank have to do about it?
\end{exercise}

\begin{exercise}[$\star\star$]\label{exo:m2:central-banks-and-the-short-rate:6}
A legacy USD~50 million loan paid three-month dollar LIBOR plus 150 basis
points and fell back to compounded SOFR. Over a 91-day period SOFR compounds
to 3.80\%. Give the all-in rate and the interest (ACT/360).
\end{exercise}

\begin{exercise}[$\star\star\star$]\label{exo:m2:central-banks-and-the-short-rate:7}
\emph{Coding.} With \texttt{firm\_rfr} and the fixings of the tutorial,
compute the September rate with a two-day lockout and with a five-day
observation shift. Give each difference from the plain compounded rate in
basis points, and explain why the shift gives exactly the lookback's number.
\end{exercise}

\begin{exercise}[$\star\star\star$]\label{exo:m2:central-banks-and-the-short-rate:8}
\emph{Find the flaw.} ``To raise the federal funds rate by a quarter point,
the Fed's trading desk sells Treasuries every morning until the rate has
risen enough.'' Correct the statement for the system the Fed ran in 2026.
In which system would it have been roughly right?
\end{exercise}

\section{Problem: The Price of a Reserve}

\begin{problem}[{Weekend problem --- the month the corridor moved}]\label{pb:m2:central-banks-and-the-short-rate:1}
A central bank runs a corridor with a deposit facility at 2.25\% and a
lending facility at 2.65\% (basis 360). Halfway through a 30-day period it
moves both by 25 basis points, to 2.50\% and 2.90\%. Bank A holds EUR~2
billion of \omterm{def:m2:central-banks-and-the-short-rate:reserves}{reserves} more than it needs over the whole period; bank B is EUR~2
billion short over the whole period. There are no other banks.

\medskip\noindent\textbf{Part I --- The facilities.}
\begin{enumerate}
\item A leaves its excess at the deposit facility for 30 days. What does it
earn?
\item B covers its shortfall at the lending facility for 30 days. What does
it pay?
\item What is the difference, and to whom does it go?
\item Why does a central bank leave a gap between the two rates at all?
\item Both rates moved by 25 basis points. Why does that move the market
rate by 25 basis points, whatever the supply of \omterm{def:m2:central-banks-and-the-short-rate:reserves}{reserves}?
\end{enumerate}

\medskip\noindent\textbf{Part II --- The interbank market.}
\begin{enumerate}[resume]
\item A lends its excess to B directly at the middle of the corridor. What
interest changes hands over the period?
\item Give A's gain over the deposit facility and B's saving over the lending
facility.
\item Nothing forces the middle. What decides where in the corridor the rate
settles?
\item A asks 5 basis points above the middle for lending unsecured to B.
Give the two new gains.
\item B could borrow from the central bank only against collateral; A lends
without. Which rate does an unsecured overnight benchmark measure, and which
does a secured one?
\end{enumerate}

\medskip\noindent\textbf{Part III --- The flood.}
\begin{enumerate}[resume]
\item The central bank buys EUR~50 billion of bonds. Where does the
interbank rate go, and why?
\item What is A's gain from lending to B now?
\item The central bank wants to raise rates by 25 basis points. What does it
change, and what does it no longer need to do?
\item Does it still need \omterm{def:m2:central-banks-and-the-short-rate:omo}{open-market operations} day to day? What is its
balance sheet for now?
\item A money-market fund with EUR~1 billion cannot use the deposit
facility. It lends to A at 6 basis points below the deposit rate for 30
days, and A deposits the cash. What does A earn on the trade?
\end{enumerate}

\medskip\noindent\textbf{Part IV --- Judgement.}
\begin{enumerate}[resume]
\item Give two reasons a central bank might prefer a corridor with scarce
\omterm{def:m2:central-banks-and-the-short-rate:reserves}{reserves}.
\item Give two reasons it might prefer a floor.
\item What does the ECB's 15-basis-point gap between its refinancing and
deposit rates since 2024 tell you about its choice?
\item State the \emph{named result}: the gain from trade in \omterm{def:m2:central-banks-and-the-short-rate:reserves}{reserves} over the
period, as a formula and in euros, and what becomes of it in a floor.
\item In one sentence: what sets the overnight rate?
\end{enumerate}
\end{problem}

\section{Interview questions}

\begin{interviewq}[$\star$ \iqroles{trader, bank}]\label{iq:m2:central-banks-and-the-short-rate:1}
What is the difference between SOFR and dollar LIBOR, and why is a SOFR
loan's margin higher than the LIBOR loan it replaced?
\end{interviewq}

\begin{interviewq}[$\star$ \iqroles{trader, researcher}]\label{iq:m2:central-banks-and-the-short-rate:2}
How can an overnight benchmark print below the central bank's deposit rate?
\end{interviewq}

\begin{interviewq}[$\star\star$ \iqroles{trader, researcher, bank}]\label{iq:m2:central-banks-and-the-short-rate:3}
Walk me through the balance sheets when a central bank does \omterm{def:m2:central-banks-and-the-short-rate:qe}{quantitative
easing}. Does it give banks more to lend?
\end{interviewq}

\begin{interviewq}[$\star\star$ \iqroles{developer, bank}]\label{iq:m2:central-banks-and-the-short-rate:4}
Why do loans that pay compounded SOFR use a lookback or an observation
shift? What is the difference between the two?
\end{interviewq}

\begin{interviewq}[$\star\star$ \iqroles{researcher, trader}]\label{iq:m2:central-banks-and-the-short-rate:5}
From market data alone, how would you tell whether a central bank runs a
floor or a \omterm{def:m2:central-banks-and-the-short-rate:corridor}{corridor system}?
\end{interviewq}

\begin{interviewq}[$\star\star\star$ \iqroles{developer}]\label{iq:m2:central-banks-and-the-short-rate:6}
You are writing the function that computes interest on a compounded-SOFR
loan. List the edge cases your tests must cover.
\end{interviewq}
